\SourceLeaf{130}
Man sieht leicht, wie sich jede Zahl $A$ in die obige Form bringen lässt. $\mu$ und $a$ ergeben sich sofort. Setzt man dann $A = 10^{\mu} a B$, so ist $b$ die Ziffer an der ersten Dezimalstelle von $B$. Setzt man ferner $B = \left(1+\frac{1}{10} b\right) C$, so ist $c$ die Ziffer an der zweiten Dezimalstelle von $C$ usw. Bei näherer Betrachtung zeigt sich, dass sich $\frac{1}{A}$ noch bequemer in die obige Form bringen lässt, ohne diesen Quotienten tatsächlich zu berechnen. Man setzt dazu nacheinander $10^{\mu} a A = B$, $\left(1+\frac{1}{10} b\right) B = C$ usw., sodass man $B$ aus $A$, $C$ aus $B$ usw. durch Multiplikation erhält, während bei der ersten Methode eine Division nötig ist. \textsc{Leonelli} hat die Rechnung bei beiden Methoden jedoch weiter verkürzt, indem er sie auf ein festes Verfahren zurückgeführt hat. Wer diese Tafeln häufig verwenden möchte, tut gut daran, sich damit vertraut zu machen. Nach einiger Übung wird es gewiss bald gelingen, die nötige Aufmerksamkeit ganz automatisch aufzubringen. Daher halten wir mit \textsc{Leonelli} den Einwand \textsc{Delambres} für voreilig, dass eine ermüdende Aufmerksamkeit nötig sei, wenn man nicht jedes Mal all die überflüssigen Nullen aufschreibt.

Man sieht leicht, dass man auf ähnliche Weise auch $A$ oder $\frac{1}{A}$ in die Form
\[
10^{\mu} a \left(1+\frac{1}{100} b\right)\left(1+\frac{1}{10000} c\right)\ldots
\]
bringen kann, wobei $a$, $b$, $c\ldots$ lauter ganze Zahlen unter 100 bezeichnen. In jedem einzelnen Fall braucht man dann nur halb so viele Glieder. Die Hilfstafel enthält zwar ebenfalls nur halb so viele Ordnungen, aber in jeder 99 Logarithmen, und muss deshalb fast fünfmal so umfangreich sein. \textsc{Leonelli} hat auch eine solche Hilfstafel beigefügt, allerdings nur für die briggischen Logarithmen und nur auf 15 Dezimalstellen. Sie enthält 485 Logarithmen. Es ist schade, dass nicht auch diese Tafel auf 20 Dezimalstellen berechnet ist.

Für die umgekehrte Aufgabe, zu einem gegebenen Logarithmus die Zahl zu finden, eignen sich \textsc{Leonellis} Tafeln ebenso gut. Zerlegt man den gegebenen Logarithmus in seine in den Tafeln enthaltenen Bestandteile, so erhält man die Folge der Faktoren der Zahl. Durch deren Multiplikation, für die sich leicht ein bequemes Verfahren finden lässt, erhält man die Zahl selbst.

